\section{Self-improvement behind one interface}
\label{sec:methods}

Every method is a class with one constructor and three calls:
\begin{lstlisting}[language=Python]
m = rsi.Reflexion(library=True)
m.fit(generate, tasks.dev)
ep = m.solve(generate, task, seed=11)
ep.accepted, ep.code, ep.usage
m.state()
\end{lstlisting}
\texttt{state()} returns what the method learned and is hashed into every episode record;
\texttt{fit()} does nothing for a method that does not learn. \texttt{rsi.benchmark()} fits each method once on the
development tasks and runs it on every held-out task and seed.

\begin{figure}[t]
\centering
% One episode: what a method sees (check() results only) and what it never sees (the
% hidden grader). Single column. Needs the TikZ libraries shadows.blur and shadings;
% pictograms come from figures/icons.tex.
\begin{tikzpicture}[
  font=\scriptsize,
  >={Stealth[length=4pt, width=3.8pt, round]},
  flow/.style={->, line width=0.9pt, line cap=round, draw=#1},
  flow/.default=fingray!85!black,
  lab/.style={font=\tiny, text=fingray!95!black, align=center, inner sep=1pt},
  title/.style={font=\footnotesize\bfseries, anchor=south west, inner sep=0pt},
  ball/.style={circle, shading=ball, ball color=#1, text=white, font=\tiny\bfseries,
               inner sep=0pt, minimum size=9pt},
  ball/.default=finblue,
  iso/.style={x={(0.95cm,-0.24cm)}, y={(0.62cm,0.32cm)}, z={(0cm,1cm)}},
  callout/.style={align=left, anchor=west, inner sep=2.2pt, text width=2.9cm, font=\tiny,
                  text=black!85, fill=white, rounded corners=2.5pt,
                  blur shadow={shadow blur steps=6, shadow xshift=0.4pt,
                               shadow yshift=-0.8pt, shadow opacity=22}},
]
% ================================================================== (b) one episode

\def\pw{6.9}
% \card{name}{x,y}{colour}{HEADER}{illustration}{caption}
\newcommand{\card}[6]{%
  \node[fill=white, rounded corners=3pt, minimum width=1.92cm, minimum height=1.26cm,
        blur shadow={shadow blur steps=7, shadow xshift=0.5pt, shadow yshift=-1pt,
                     shadow opacity=28}] (#1) at (#2) {};
  \begin{scope}
    \clip[rounded corners=3pt] (#1.south west) rectangle (#1.north east);
    \shade[left color=#3, right color=#3!75!white] ($(#1.north west) + (0, -0.27)$) rectangle (#1.north east);
  \end{scope}
  \draw[#3!55, line width=0.4pt, rounded corners=3pt] (#1.south west) rectangle (#1.north east);
  \node[font=\tiny\bfseries, text=white, anchor=center] at ($(#1.north) + (0, -0.135)$) {#4};
  \begin{scope}[shift={($(#1.center) + (0, -0.06)$)}] #5 \end{scope}
  \node[font=\tiny, text=black!80, anchor=south] at ($(#1.south) + (0, 0.02)$) {#6};}



\shade[top color=finblue!10, bottom color=finblue!3, rounded corners=7pt, draw=finblue!35,
       line width=0.5pt] (0, 4.27) rectangle (\pw, 0.95);
\shade[top color=fingray!17, bottom color=fingray!6, rounded corners=7pt, draw=fingray!45,
       line width=0.5pt] (0, 0.59) rectangle (\pw, -1.06);
\node[lab, anchor=south east, font=\tiny\bfseries, text=finblue!75!black] at (\pw-0.12, 1.02)
  {VISIBLE TO THE METHOD};
\node[lab, anchor=south west, font=\tiny\bfseries, text=fingray!95!black] at (0.12, -1.0)
  {HIDDEN FROM METHODS};

\card{method}{1.15, 3.5}{finblue}{METHOD}{%
  \fill[fingray!85!black, rounded corners=1pt] (-0.55, -0.13) rectangle (0.55, 0.2);
  \foreach \i/\c/\l in {0/finorange/0.5, 1/finaqua/0.75, 2/finblue!60/0.4}
    \fill[\c] (-0.45, {0.12-\i*0.09}) rectangle ({-0.45+\l}, {0.085-\i*0.09});}{\texttt{solve()}}
\card{meter}{3.45, 3.5}{finblue}{METER}{%
  \foreach \a/\b/\c in {180/120/finaqua, 120/60/finorange!70!yellow, 60/0/finorange}
    \draw[\c, line width=2.2pt] (\a:0.3) arc[start angle=\a, end angle=\b, radius=0.3];
  \draw[fingray!90!black, line width=0.8pt, line cap=round] (0,0) -- (50:0.26);
  \fill[fingray!90!black] (0,0) circle (0.045);}{equal budget}
\card{llm}{5.75, 3.5}{fingray}{YOUR MODEL}{%
  \node[circle, shading=ball, ball color=fingray!65, minimum size=0.46cm] at (0, 0.02) {};
  \pic[draw=white, fill=white, line width=0.6pt] at (0, 0.02) {fin robot};}{\texttt{generate()}}
\card{judge}{1.15, 1.66}{finorange}{JUDGE}{%
  \foreach \yy/\c/\t in {0.17/finaqua/finite, 0.06/finaqua/split, -0.05/finorange/lag} {
    \fill[\c] (-0.55, \yy) circle (0.045);
    \node[font=\fontsize{4}{4}\selectfont\ttfamily, anchor=west, inner sep=0pt] at (-0.47, \yy) {\t};
    \fill[fingray!25] (0.05, {\yy-0.02}) rectangle (0.55, {\yy+0.02});}}{\texttt{check()}}
\card{grader}{3.62, -0.22}{fingray!90!black}{HIDDEN GRADER}{%
  \begin{scope}[iso, scale=0.32, shift={(-0.6, -0.4, -0.55)}]
    \fill[fingray!80!black] (1,0,0) -- (1,1,0) -- (1,1,1) -- (1,0,1) -- cycle;
    \fill[fingray!60!black] (0,0,0) -- (1,0,0) -- (1,0,1) -- (0,0,1) -- cycle;
    \fill[fingray!35] (0,0,1) -- (1,0,1) -- (1,1,1) -- (0,1,1) -- cycle;
  \end{scope}
  \pic[draw=white, fill=white, line width=0.5pt, scale=0.8] at (-0.07, -0.05) {fin lock};}{vs.\ reference}
\card{report}{5.9, -0.22}{finaqua}{BENCHMARK}{%
  \foreach \i/\h/\c in {0/0.18/fingray!45, 1/0.30/finblue, 2/0.24/finorange, 3/0.38/finaqua}
    \fill[\c, rounded corners=0.6pt] ({-0.42+\i*0.23}, -0.18) rectangle ({-0.26+\i*0.23}, {-0.18+\h});}{valid, cost}

% the barrier
\draw[fingray!90!black, line width=1pt, dash pattern=on 4pt off 2.5pt] (0.1, 0.77) -- (\pw-0.1, 0.77);
\node[ball=fingray!80, minimum size=0.46cm] (barrier) at (\pw - 0.5, 0.77) {};
\pic[draw=white, fill=white, line width=0.55pt] at (barrier) {fin lock};
\node[lab, anchor=east, font=\tiny\itshape, fill=white, rounded corners=1.5pt, inner sep=1.5pt]
  at ($(barrier.west) + (-0.06, 0)$) {no grade crosses this line};

% arrows and steps
\draw[flow=finblue!80!black] ([yshift=7pt]method.east) -- ([yshift=7pt]meter.west);
\draw[flow=finblue!80!black] ([yshift=7pt]meter.east) -- ([yshift=7pt]llm.west);
\draw[flow=finblue!80!black] ([yshift=-7pt]llm.west) -- ([yshift=-7pt]meter.east);
\draw[flow=finblue!80!black] ([yshift=-7pt]meter.west) -- ([yshift=-7pt]method.east);
\node[ball] at ($([yshift=7pt]meter.east)!0.5!([yshift=7pt]llm.west) + (0, 0.27)$) {1};
\node[ball] at ($([yshift=-7pt]llm.west)!0.5!([yshift=-7pt]meter.east) + (0, -0.27)$) {2};
\draw[flow=finblue!80!black] ([xshift=-10pt]method.south) -- ([xshift=-10pt]judge.north);
\draw[flow=finorange!85!black] ([xshift=10pt]judge.north) -- ([xshift=10pt]method.south);
\node[ball] at ($([xshift=-10pt]method.south)!0.5!([xshift=-10pt]judge.north) + (-0.25, 0)$) {3};
\node[ball=finorange] at ($([xshift=10pt]judge.north)!0.5!([xshift=10pt]method.south) + (0.25, 0)$) {4};
\node[lab, anchor=west, align=left] at ($(judge.east) + (0.12, 0.05)$)
  {ran?\ error,\\failed checks\\and their causes};
\draw[flow=fingray!95!black] (judge.south) |- (grader.west);
\node[ball=fingray!85] at ($(judge.south |- grader.west) + (-0.25, 0.3)$) {5};
\node[lab, anchor=south west, align=left] at ($(judge.south |- grader.west) + (0.1, 0.03)$)
  {final\\program};
\draw[flow=fingray!95!black] (grader.east) -- (report.west);
\node[ball=fingray!85] at ($(grader.east)!0.5!(report.west) + (0, 0.27)$) {6};


\end{tikzpicture}

\caption{One episode. (1)~Each prompt and (2)~each program pass through the meter; (3)~the
method sends code to the judge and (4)~sees only what \texttt{check()} returns. (5)~After
\texttt{solve()} returns, the hidden grader runs on the final program, and (6)~its verdict
goes only to the benchmark report.}
\label{fig:episode}
\end{figure}

\paragraph{What a method sees.}
\texttt{solve()} gets the task text and a judge on which it can only call
\texttt{check(code)}. That returns whether the program ran, its error, each check's result
and the causes of the failed ones, never the grade (Figure~\ref{fig:episode}). Learning
methods may use development grades to keep or reject a candidate, but no text a model reads
contains a grade.

\paragraph{Equal budgets.}
A meter refuses any call beyond the budget and records which limit stopped the episode.
Each episode gets \BudgetAttempts{} program executions, \BudgetCalls{} model calls
(reflections and reviews included) and \BudgetTokens{} generated tokens. Learning methods
also get \FitEpisodes{} development episodes, \FitProposals{} proposals, \FitCalls{} calls
and \FitTokens{} tokens.

\paragraph{Methods.}
The \NMethods{} methods (Appendix Table~\ref{tab:methods}) keep the mechanism each is known
for: \texttt{NoImprove} (one program); \texttt{SelfDebug}, which returns only errors
\citep{chen2024selfdebug}; \texttt{CheckFeedback}, which returns the failed checks with
causes and cards; \texttt{Reflexion}, which adds a written reflection before each retry
\citep{shinn2023reflexion}; \texttt{Team}, with one reviewer per rule family;
\texttt{Evolve}, which evolves a prompt addendum on development tasks
\citep{agrawal2026gepa}; and \texttt{PlaybookLearning}
\citep{zhao2024expel,suzgun2026dynamic}. The playbook proposer sees development traces
without grades, and one proposal edits one entry. A candidate is accepted only if its score
does not drop on either of two disjoint halves of the development tasks, rises on one, and
no task that passed every check before now fails one (the \emph{no-regression} rule;
Appendix Figure~\ref{fig:gate}).

\paragraph{Skill arms.}
Two switches split the skills: \texttt{cards} adds the task's rule cards to the prompt and
the feedback; \texttt{library} adds the interface text and allows
\texttt{import fin\_skills}. Their \NArms{} combinations (\texttt{none}, \texttt{cards},
\texttt{library}, \texttt{cards+library}) all keep the task definition and see the failed
checks with their causes. Within an arm, \texttt{NoImprove} versus \texttt{CheckFeedback}
measures the checks and \texttt{CheckFeedback} versus \texttt{PlaybookLearning} measures
self-improvement; for one method, comparing arms measures the skills (Appendix
Figure~\ref{fig:arms}).
